\documentclass{article}
\usepackage{graphicx}
\usepackage{amsmath}
\usepackage{amssymb}
\usepackage{algorithm}
\usepackage{algpseudocode}
\usepackage{booktabs}
\usepackage{hyperref}
\usepackage{xcolor}
\usepackage{listings}
\usepackage{geometry}
\geometry{margin=2.5cm}

\lstset{
  basicstyle=\ttfamily\small,
  breaklines=true,
  frame=single,
  numbers=left,
  numberstyle=\tiny,
  keywordstyle=\color{blue},
  commentstyle=\color{gray},
  stringstyle=\color{red}
}

\title{Efficient Methods for Solving the Traveling Salesman Problem}
\author{Filippo Tiberio \and Herman Serain}
\date{March 2026}

\begin{document}

\maketitle

\begin{abstract}
The Traveling Salesman Problem (TSP) is one of the most well-known and widely studied problems
in combinatorial optimization. Despite its simple formulation, the problem quickly becomes
computationally intractable as the number of cities increases. Because of its theoretical
importance and its wide range of practical applications, the TSP has been the subject of
extensive research for several decades.

This report introduces the Traveling Salesman Problem and discusses the main techniques used
to solve it efficiently. After presenting the formal definition of the problem and its
mathematical representation, we briefly review the computational challenges associated with
it. The goal of this work is to provide an overview of the main algorithmic approaches used
to address the TSP, highlighting the trade-offs between optimality and computational
efficiency. In particular, we examine construction and refinement heuristics, metaheuristic
strategies, and exact methods based on integer linear programming.
\end{abstract}

\tableofcontents
\newpage

% ============================================================
\section{Introduction}
% ============================================================

The Traveling Salesman Problem (TSP) is a classical problem in combinatorial optimization
and theoretical computer science. Informally, the problem can be described as follows: given
a set of cities and the distances between every pair of them, the objective is to determine
the shortest possible route that visits each city exactly once and returns to the starting
city.

Although the problem is easy to state and understand, finding the optimal solution is
computationally challenging. The difficulty arises from the extremely large number of
possible tours that must be considered. For a problem with $n$ cities, the number of
potential routes grows factorially with $n$. Even for relatively small instances, the number
of possible solutions becomes prohibitively large, making exhaustive search impractical.

The TSP is classified as an NP-hard problem in computational complexity theory. This means
that no algorithm is currently known that can solve all instances of the problem in
polynomial time. As a consequence, a large portion of research has focused on developing
more efficient algorithms and approximation techniques that can produce optimal or
near-optimal solutions within a reasonable computational time.

Beyond its theoretical importance, the Traveling Salesman Problem has numerous practical
applications. It appears naturally in many real-world scenarios where a set of locations must
be visited in the most efficient order. Examples include logistics and delivery route
planning, vehicle routing problems, manufacturing processes such as printed circuit board
drilling, and even certain problems in computational biology.

Because of its importance, the TSP has become a benchmark problem used to test and compare
optimization algorithms. Over the years, researchers have developed a wide variety of
solution approaches, including exact algorithms, approximation algorithms, and heuristic
methods. Each of these approaches presents different advantages in terms of accuracy,
computational cost, and scalability.

In this report, we introduce the fundamental concepts behind the Traveling Salesman Problem
and discuss the main approaches used to solve it efficiently. We begin by presenting the
formal definition of the problem and its mathematical representation, which will serve as the
foundation for the methods discussed later.

% ============================================================
\section{Background Material}
% ============================================================

\subsection{What is the Traveling Salesman Problem}

The Traveling Salesman Problem is a combinatorial optimization problem defined on a set of
nodes representing cities and edges representing the distances between them. The objective
is to determine the shortest closed tour that visits every city exactly once and returns to
the starting point.

From a mathematical perspective, the problem is typically modeled using a weighted graph.
Let $G = (V, E)$ be a graph where $V$ is the set of vertices representing the cities and $E$
is the set of edges connecting pairs of cities. Each edge $(i,j)$ is associated with a cost
or distance $d_{ij}$ representing the travel cost between city $i$ and city $j$.

The goal is to find a tour that minimizes the total travel distance while ensuring that each
city is visited exactly once. This can be expressed as the minimization of the total cost of
the tour:

\[
\min \sum_{i=1}^{n} d_{\pi(i),\pi(i+1)}
\]

where $\pi$ represents a permutation of the vertices and $\pi(n+1) = \pi(1)$ ensures that
the route returns to the starting city.

The TSP can appear in several different variants depending on the properties of the distance
matrix. One of the most common forms is the \textit{symmetric Traveling Salesman Problem},
in which the distance between two cities is the same in both directions:

\[
d_{ij} = d_{ji}
\]

In contrast, the \textit{asymmetric Traveling Salesman Problem} allows the distance from
city $i$ to city $j$ to differ from the distance from city $j$ to city $i$. This situation
may arise in real-world applications such as road networks with one-way streets or different
travel costs depending on direction.

A key difficulty of the TSP lies in the rapid growth of the solution space. For a symmetric
TSP with $n$ cities, the number of possible tours is approximately

\[
\frac{(n-1)!}{2}
\]

since many tours are equivalent due to symmetry and rotation. Even for moderately sized
problems, the number of possible tours becomes extremely large. For example, with only $20$
cities there are already more than $10^{18}$ possible tours.

Because of this combinatorial explosion, solving the TSP efficiently requires sophisticated
optimization techniques. Over the years, researchers have developed a wide range of
approaches, including exact algorithms capable of guaranteeing optimal solutions, as well as
heuristic and metaheuristic methods that can find high-quality solutions for very large
instances.

Understanding the structure of the TSP and the challenges associated with it is essential
before exploring these solution methods. In the following sections, we will discuss some of
the most important algorithmic strategies developed to tackle this problem.

% -----------------------------------------------------------
\subsection{IBM CPLEX Optimizer}
% -----------------------------------------------------------

IBM CPLEX Optimizer, hereafter referred to simply as CPLEX, is a high-performance
mathematical programming solver developed and maintained by IBM. It is widely used in
operations research and combinatorial optimization to solve a broad class of problems,
including linear programs (LP), mixed-integer programs (MIP), and quadratic programs (QP).
For the purposes of this work, CPLEX is employed to solve the TSP formulated as an integer
linear program (ILP), providing provably optimal solutions for moderate-sized instances.

% CPLEX operates through a branch-and-cut framework for integer programs. Starting from the
% LP relaxation of the problem, it explores a search tree where at each node the relaxation is
% solved and, if the solution is fractional, branching variables are selected and cutting
% planes are added to tighten the formulation. This combination of branch-and-bound with
% cutting plane methods allows CPLEX to solve instances of significant size that would be
% intractable by pure enumeration.

CPLEX exposes a rich C API, through
which the user can programmatically build a model by adding variables, objective functions,
and constraints, and then invoke the solver. The key API objects are the \texttt{CPXENVptr}
environment handle, which manages solver parameters and licensing, and the \texttt{CPXLPptr}
problem handle, which stores the mathematical model.

% \paragraph{Installation.}
% CPLEX is available as part of the IBM Decision Optimization suite. Academic users can obtain
% a free full-featured license through the IBM Academic Initiative program. The installation
% procedure is outlined below at a high level:

% \begin{enumerate}
%     \item Register on the IBM Academic Initiative portal and download the CPLEX installer for
%           the target operating system (Windows, Linux, or macOS).
%     \item Run the installer and follow the guided setup wizard, selecting the desired
%           installation directory.
%     \item Set the required environment variables: \texttt{CPLEX\_HOME} should point to the
%           CPLEX installation directory, and the \texttt{PATH} (or \texttt{LD\_LIBRARY\_PATH}
%           on Linux) should include the path to the CPLEX shared libraries.
%     \item When compiling a C project that uses CPLEX, link against the \texttt{cplex} library
%           and include the header \texttt{ilcplex/cplex.h}. A typical compilation command on
%           Linux resembles:
% \begin{lstlisting}[language=bash]
% gcc -O2 -o tsp_solver main.c tsp.c \
%     -I$CPLEX_HOME/include \
%     -L$CPLEX_HOME/lib/x86-64_linux/static_pic \
%     -lcplex -lpthread -lm -ldl
% \end{lstlisting}
%     \item Verify the installation by running the bundled example programs provided in the
%           CPLEX \texttt{examples/} directory.
% \end{enumerate}

For this project, CPLEX version 22.1 was used on a Linux x86-64 machine, a WIN11 machine and a macOS Tahoe (version 26.5) machine. The solver is invoked exclusively through its C callable library.

% -----------------------------------------------------------
\subsection{Performance Profiles}
% -----------------------------------------------------------

When comparing multiple optimization algorithms on a set of benchmark instances, a common
challenge is that no single scalar metric captures the full picture of an algorithm's
behavior. An algorithm may be superior on some instances and inferior on others; moreover,
the degree of superiority matters. The \textit{performance profile}, introduced by Dolan and
Moré \cite{dolan2002benchmarking}, is a widely adopted tool in the operations research
community for conducting such comparisons in a statistically robust and visually informative
manner.

\paragraph{Definition.}
Let $\mathcal{S}$ be the set of solvers (algorithms) to be compared and $\mathcal{P}$ be the
set of benchmark problem instances. For each solver $s \in \mathcal{S}$ and each instance
$p \in \mathcal{P}$, let $t_{p,s}$ denote the performance metric of interest (e.g.,
solution cost or CPU time). The \textit{performance ratio} of solver $s$ on instance $p$ is
defined as:

\[
r_{p,s} = \frac{t_{p,s}}{\min_{s' \in \mathcal{S}} t_{p,s'}}
\]

This ratio measures how far solver $s$ is from the best solver on instance $p$. A ratio of
$1$ means the solver achieved the best result on that instance. The performance profile of
solver $s$ is then the cumulative distribution function of the performance ratios:

\[
\rho_s(\tau) = \frac{1}{|\mathcal{P}|} \left| \left\{ p \in \mathcal{P} : r_{p,s} \leq \tau \right\} \right|
\]

where $\tau \geq 1$ is the performance ratio threshold. Intuitively, $\rho_s(\tau)$ gives
the fraction of instances on which solver $s$ is within a factor $\tau$ of the best solver.

\paragraph{Interpretation.}
The value $\rho_s(1)$ represents the fraction of instances on which solver $s$ achieves the
best performance among all solvers — a measure of \textit{winning frequency}. As $\tau$
increases, $\rho_s(\tau)$ approaches $1$ (the solver eventually becomes competitive on all
instances). An algorithm with a profile that lies above and to the left of all others is
uniformly superior. When profiles cross, the interpretation depends on the specific operating
region of $\tau$ that is of practical interest.

In this work, performance profiles are used to compare the solution quality achieved by the
various heuristic, metaheuristic and math-heuristic algorithms described in the following sections, evaluated
over a suite of randomly generated TSP benchmark instances. The metric adopted is the normalized solution cost,
defined as the ratio of the cost obtained by each method to the best known solution cost for
each instance. This normalization makes performance profiles scale-independent and suitable
for aggregating results across instances of different sizes.


% ============================================================
\section{Heuristic Methods}
% ============================================================

Heuristic methods do not guarantee optimality but are designed to find good-quality solutions
within a reasonable amount of time. They are particularly valuable for large TSP instances
where exact methods become computationally prohibitive. In this work, heuristics are
organized into two categories: \textit{construction heuristics}, which build a complete tour
from scratch, and \textit{refinement heuristics}, which iteratively improve an existing tour.

% -----------------------------------------------------------
\subsection{Construction Heuristics}
% -----------------------------------------------------------

Construction heuristics produce an initial feasible solution by incrementally assembling a
tour. While the resulting solutions are generally not optimal, they serve as high-quality
starting points for subsequent refinement procedures.

\subsubsection{Nearest Neighbor Heuristic}

The Nearest Neighbor (NN) heuristic is one of the simplest and most widely known greedy
construction methods for the TSP. Its appeal lies in its conceptual simplicity and low
computational overhead. The algorithm builds a tour by iteratively extending the current
partial tour with the unvisited city that is closest to the last city added.

\paragraph{Algorithm Description.}
Starting from an initial city, the heuristic greedily selects at each step the nearest
unvisited city and appends it to the current tour. Once all cities have been visited, the
tour is closed by returning to the starting city. The procedure is repeated from every city
as a starting node, and the best tour found across all starting points is retained as the
final solution.

The pseudocode for a single iteration of the Nearest Neighbor heuristic is given in
Algorithm~\ref{alg:nn}.

\begin{algorithm}[H]
\caption{Nearest Neighbor Heuristic (single starting node)}
\label{alg:nn}
\begin{algorithmic}[1]
\Require Instance with $n$ cities, starting node $s$
\Ensure Tour $\pi$ and its cost $C$
\State $\pi \leftarrow [0, 1, \dots, n{-}1]$
\State \textbf{swap} $\pi[0] \leftrightarrow \pi[s]$
\State $C \leftarrow 0$
\For{$i = 0$ \textbf{to} $n-3$}
    \State $\text{best} \leftarrow \arg\min_{j \in \{i+1,\dots,n-1\}} d_{\pi[i],\,\pi[j]}$
    \State \textbf{swap} $\pi[i{+}1] \leftrightarrow \pi[\text{best}]$
    \State $C \leftarrow C + d_{\pi[i],\,\pi[i+1]}$
\EndFor
\State $C \leftarrow C + d_{\pi[n-2],\,\pi[n-1]} + d_{\pi[n-1],\,\pi[0]}$ \Comment{Last edge + closing edge}
\State \Return $\pi$, $C$
\end{algorithmic}
\end{algorithm}

\paragraph{Implementation Details.}
In our implementation, the distances between all pairs of cities are precomputed and stored
in a compact upper triangular matrix to avoid redundant calculations during the heuristic
loop. Access to the precomputed distance between nodes $i$ and $j$ (with $i < j$) is
realized through the index formula:

\[
\text{index}(i, j, N) = i \cdot N - \frac{i(i+1)}{2} + (j - i - 1)
\]

which maps the upper triangle of the distance matrix to a one-dimensional array of size
$\binom{N}{2}$. The multi-start strategy, running the heuristic from each city as a
starting point, is implemented with a randomized ordering of starting nodes to allow
early termination within a given time budget while still exploring diverse starting points.

\paragraph{Complexity.}
Each iteration of the Nearest Neighbor heuristic from a fixed starting node requires $O(n^2)$
time in the basic implementation, as for each of the $n$ cities to be placed, all remaining
unvisited cities must be examined. Running the multi-start variant over all $n$ starting
nodes results in an overall complexity of $O(n^3)$.

% \paragraph{Quality of Solutions.}
% Despite its simplicity, the Nearest Neighbor heuristic typically produces solutions that are
% within 20--25\% of the optimal tour length for Euclidean instances. However, it is known to
% suffer from the \textit{end-of-tour} effect: the final cities added to the tour are often
% placed sub-optimally because the greedy choices made early on constrain later options.
% Multi-start mitigates this issue to some extent.

\subsubsection{Extra Mileage Heuristic}

The Extra Mileage (EM) heuristic, also known as the cheapest insertion heuristic, builds
a tour by starting with a minimal sub-tour and incrementally inserting the remaining cities
at the position that incurs the least additional cost. Compared to the Nearest Neighbor
approach, it tends to produce more balanced tours because it considers the impact of each
insertion on the existing partial tour.

\paragraph{Algorithm Description.}
The heuristic begins by identifying the two most distant cities, which form the initial
sub-tour of length 2 (a back-and-forth path between those two cities). At each subsequent
step, for every unvisited city $k$ and every edge $(i, j)$ in the current partial tour,
the \textit{extra mileage} required to insert $k$ between $i$ and $j$ is computed as:

\[
\Delta(k, i, j) = d_{ik} + d_{kj} - d_{ij}
\]

The city $k^*$ and insertion position $(i^*, j^*)$ that minimize $\Delta$ are selected, and
$k^*$ is inserted between $i^*$ and $j^*$ in the tour. This process repeats until all cities
have been inserted.

\begin{algorithm}[H]
\caption{Extra Mileage Heuristic}
\label{alg:em}
\begin{algorithmic}[1]
\Require Instance with $n$ cities
\Ensure Tour $\pi$ and its cost $C$
\State $(u, v) \leftarrow \arg\max_{i < j}\; d_{ij}$
\State $\pi \leftarrow [u,\, v]$, $\text{visited}[u] \leftarrow \text{visited}[v] \leftarrow \text{true}$, $m \leftarrow 2$
\While{$m < n$}
    \State $\Delta^* \leftarrow +\infty$, $k^* \leftarrow -1$, $p^* \leftarrow -1$
    \For{each unvisited city $k$}
        \For{$i = 0$ \textbf{to} $m - 1$}
            \State $j \leftarrow (i + 1) \bmod m$
            \State $\Delta \leftarrow d_{\pi[i],\,k} + d_{k,\,\pi[j]} - d_{\pi[i],\,\pi[j]}$
            \If{$\Delta < \Delta^*$}
                \State $\Delta^* \leftarrow \Delta$, $k^* \leftarrow k$, $p^* \leftarrow i + 1$
            \EndIf
        \EndFor
    \EndFor
    \State Shift right $\pi[p^*\,..\,m]$ by one position \Comment{Make room}
    \State $\pi[p^*] \leftarrow k^*$, $\text{visited}[k^*] \leftarrow \text{true}$, $m \leftarrow m + 1$
\EndWhile
\State $C \leftarrow \textsc{TourCost}(\pi)$
\State \Return $\pi$, $C$
\end{algorithmic}
\end{algorithm}

\paragraph{Complexity.}
The Extra Mileage heuristic has a worst-case time complexity of $O(n^3)$, since at each of
the $n-2$ insertion steps, it must evaluate up to $n$ unvisited cities against up to $n$
edges in the current partial tour. In practice, the constant factors are slightly larger than
those of the Nearest Neighbor heuristic, making EM somewhat slower for very large instances.

% \paragraph{Quality of Solutions.}
% The Extra Mileage heuristic generally produces higher-quality initial solutions than the
% Nearest Neighbor approach, as the cheapest insertion criterion tends to avoid the severe
% detours that the greedy nearest-neighbor strategy can introduce at the end of the tour. It
% is, however, a single-start method and is typically applied once rather than in a multi-start
% fashion.

% -----------------------------------------------------------
\subsection{Refinement Heuristics}
% -----------------------------------------------------------

Refinement heuristics take an existing feasible tour and attempt to improve it by performing
local modifications. They are typically applied as a post-processing step after a
construction heuristic has produced an initial solution, and they continue iterating until
no further improvement is possible or a time budget is exhausted.

\subsubsection{2-opt Refinement}

The 2-opt algorithm, originally proposed by Lin (1965) \cite{lin1965}, is the most widely
used local search refinement method for the TSP. It works by repeatedly selecting two edges
in the current tour and replacing them with two alternative edges, effectively reversing the
segment of the tour between the two selected edges. This operation is called a \textit{2-opt
move}.

\paragraph{2-opt Move.}
Given a tour $\pi = [\pi_1, \pi_2, \ldots, \pi_n]$, a 2-opt move selects two non-adjacent
edges $(\pi_i, \pi_{i+1})$ and $(\pi_j, \pi_{j+1})$ and replaces them with the edges
$(\pi_i, \pi_j)$ and $(\pi_{i+1}, \pi_{j+1})$. The segment $[\pi_{i+1}, \ldots, \pi_j]$
is then reversed, yielding a new tour. The cost change of such a move is:

\[
\Delta C = \left( d_{\pi_i, \pi_j} + d_{\pi_{i+1}, \pi_{j+1}} \right)
         - \left( d_{\pi_i, \pi_{i+1}} + d_{\pi_j, \pi_{j+1}} \right)
\]

A move is accepted if and only if $\Delta C < 0$, i.e., the new tour is strictly shorter.

\paragraph{Best-Improvement Strategy.}
Our implementation uses the \textit{best-improvement} strategy: in each iteration, all $O(n^2)$ possible 2-opt moves are evaluated, and the one
with the largest improvement ($\Delta C$ most negative) is applied. This contrasts with the
\textit{first-improvement} strategy, which applies the first improving move found. While
best-improvement typically requires more evaluations per iteration, it often converges to a better local optimum.

\begin{algorithm}[H]
\caption{2-opt Refinement}
\label{alg:2opt}
\begin{algorithmic}[1]
\Require Tour $\pi$ of $n$ cities, cost $C$
\Ensure Locally optimal tour $\pi$, cost $C$
\Repeat
    \State $\Delta^* \leftarrow 0$, $a^* \leftarrow -1$, $b^* \leftarrow -1$
    \For{$i \leftarrow 0$ \textbf{to} $n - 2$}
        \For{$j \leftarrow i + 2$ \textbf{to} $n - 1$, $j \neq (i + n - 1) \bmod n$}
            \State $i' \leftarrow (i+1) \bmod n$, $\quad j' \leftarrow (j+1) \bmod n$
            \State $\Delta C \leftarrow \bigl(d_{\pi[i],\pi[j]} + d_{\pi[i'],\pi[j']}\bigr) - \bigl(d_{\pi[i],\pi[i']} + d_{\pi[j],\pi[j']}\bigr)$
            \If{$\Delta C < \Delta^*$}
                \State $\Delta^* \leftarrow \Delta C$, $\quad a^* \leftarrow i$, $\quad b^* \leftarrow j$
            \EndIf
        \EndFor
    \EndFor
    \If{$\Delta^* < 0$}
        \State Reverse segment $\pi[a^*{+}1\,..\,b^*]$
        \State $C \leftarrow C + \Delta^*$
    \EndIf
\Until{$\Delta^* \geq 0$}
\State \Return $\pi$, $C$
\end{algorithmic}
\end{algorithm}

\paragraph{Complexity.}
Each iteration of the 2-opt algorithm requires evaluating $O(n^2)$ pairs of edges. The
number of iterations until convergence can vary widely, but in practice the algorithm
typically converges after a small polynomial number of steps. The overall worst-case
complexity is $O(n^3)$ per local optimum reached, though empirically the behavior is
substantially better.

\paragraph{2-opt Local Optima.}
A tour is called \textit{2-optimal} if no 2-opt move improves it. Although 2-optimal tours
are generally much better than the initial greedy construction, they are not guaranteed to be globally optimal.

% In practice, 2-opt applied after the
% Nearest Neighbor heuristic typically reduces the tour cost by 10--20\% on benchmark
% instances.
\begin{figure}[H]
    \centering
    \includegraphics[width=1\linewidth]{Figure_1.png}
    \caption{NN Multistart + 2 OPT (Blue) VS EM + 2 OPT (Orange)}
    \label{fig:placeholder}
\end{figure}

% ============================================================
\section{Metaheuristic Methods}
% ============================================================

Metaheuristics are high-level strategies that guide the search process beyond the limits
of pure local search. Unlike basic heuristics, which tend to get trapped in local optima,
metaheuristics employ mechanisms such as controlled randomization, memory structures, and
perturbation operators to explore a larger portion of the solution space. In this work, we
consider three metaheuristic approaches: GRASP (a construction metaheuristic), and VNS and
Tabu Search (refinement metaheuristics).

% -----------------------------------------------------------
\subsection{Greedy Randomized Adaptive Search Procedure (GRASP)}
% -----------------------------------------------------------

The Greedy Randomized Adaptive Search Procedure (GRASP), introduced by Feo and Resende
\cite{feo1995greedy}, is a multi-start metaheuristic that combines greedy construction with
randomization. Each GRASP iteration consists of two phases: a randomized construction phase
that builds a feasible solution, followed by a local search phase that refines it. The best
solution found across all iterations is returned.

\paragraph{Randomized Construction.}
In the standard Nearest Neighbor heuristic, the greedy choice is deterministic: at each step
the closest unvisited city is selected. GRASP introduces controlled randomization by
perturbing the distance used to rank candidate cities. Specifically, for each candidate edge
of true length $d$, a perturbed distance $\tilde{d}$ is computed as:

\[
\tilde{d} = d \cdot \left(1 + \alpha \cdot U(0,1) \right)
\]

where $U(0,1)$ is a uniform random variable and $\alpha \in [0, 1]$ is the \textit{GRASP
randomization parameter}. When $\alpha = 0$, the heuristic reduces to the deterministic
greedy Nearest Neighbor; as $\alpha$ increases, the selection becomes more random, allowing
the algorithm to explore a broader region of the solution space.

\begin{algorithm}[H]
\caption{GRASP Construction Phase (Randomized Nearest Neighbor)}
\label{alg:grasp_constr}
\begin{algorithmic}[1]
\Require Instance, starting node $s$, randomization parameter $\alpha$
\Ensure Tour $\pi$ and its cost $C$
\State $\pi \leftarrow [s]$, $\text{visited}[s] \leftarrow \text{true}$, $C \leftarrow 0$
\While{$|\pi| < n$}
    \State $\text{current} \leftarrow \pi[\text{last}]$
    \State $\tilde{d}^* \leftarrow +\infty$, $\text{best\_node} \leftarrow -1$, $\text{real\_dist} \leftarrow 0$
    \For{each unvisited city $j$}
        \State $\tilde{d}_j \leftarrow d_{\text{current},j} \cdot (1 + \alpha \cdot U(0,1))$
        \If{$\tilde{d}_j < \tilde{d}^*$}
            \State $\tilde{d}^* \leftarrow \tilde{d}_j$, $\text{best\_node} \leftarrow j$
            \State $\text{real\_dist} \leftarrow d_{\text{current},j}$
        \EndIf
    \EndFor
    \State Append $\text{best\_node}$ to $\pi$
    \State $\text{visited}[\text{best\_node}] \leftarrow \text{true}$
    \State $C \leftarrow C + \text{real\_dist}$
\EndWhile
\State $C \leftarrow C + d_{\pi[n], \pi[1]}$
\State \Return $\pi$, $C$
\end{algorithmic}
\end{algorithm}

\paragraph{Overall GRASP Procedure.}
The full GRASP algorithm wraps the construction phase inside a multi-start loop, applying
2-opt refinement after each construction to improve the solution. The best tour found across
all iterations is maintained as the incumbent.

\begin{algorithm}[H]
\caption{GRASP for TSP}
\label{alg:grasp}
\begin{algorithmic}[1]
\Require Instance, parameter $\alpha$, time limit $T$
\Ensure Best tour $\pi^*$, cost $C^*$
\State $C^* \leftarrow +\infty$
\For{each starting node $s$ (in random order)}
    \State $(\pi, C) \leftarrow \text{GRASP-Construct}(s, \alpha)$
    \State $(\pi, C) \leftarrow \text{2-opt}(\pi, C)$
    \If{$C < C^*$}
        \State $C^* \leftarrow C$, $\pi^* \leftarrow \pi$
    \EndIf
    \If{elapsed time $> T$} \textbf{break} \EndIf
\EndFor
\State \Return $\pi^*$, $C^*$
\end{algorithmic}
\end{algorithm}

\paragraph{Role of $\alpha$.}
The parameter $\alpha$ governs the exploration-exploitation trade-off. Small values of $\alpha$
keep the construction close to the greedy solution, yielding high-quality individual
solutions but limited diversity. Large values of $\alpha$ produce highly diverse solutions
but at the cost of individual quality. Optimal values of $\alpha$ are typically determined
empirically through parameter tuning.

\subsubsection{Hyperparameter Tuning for GRASP}

The performance of GRASP is critically dependent on the choice of the randomization parameter
$\alpha$. To systematically identify the best value, a grid search was conducted over the
range $\alpha \in \{0.0,0.05, 0.1, 0.2, 0.3, 0.4, 0.5\}$, evaluating each configuration on a
set of benchmark instances under a fixed time budget. Performance profiles (Section 2.3)
were used to compare configurations in a scale-independent manner.

% [PLACEHOLDER: INSERT PERFORMANCE PROFILE PLOT HERE]
\begin{figure}[H]
    \centering
    \includegraphics[width=1\linewidth]{GRASP_PP.png}
\end{figure}
% Results indicate that \alpha \approx 0.1 achieves the best balance between solution
% diversity and construction quality. The profile for \alpha = 0.1 dominates the profiles
% for all other values for \tau in the range [1.0, 1.15], and achieves a win rate of 
% approximately 62% of instances. Values of \alpha > 0.3 show degraded performance due 
% to excessive randomization, while \alpha = 0.0 (pure greedy) performs competitively 
% on small instances but is consistently outperformed on larger ones.

The results of this analysis are reported in Table~\ref{tab:grasp_alpha} and discussed in
terms of performance profiles in Figure~\ref{fig:grasp_profile}.

\begin{table}[H]
\centering
\caption{Average normalized solution cost for different values of $\alpha$ (GRASP). Lower
is better. Results are averages over 30 benchmark instances.}
\label{tab:grasp_alpha}
\begin{tabular}{@{}ccc@{}}
\toprule
$\alpha$ & Avg. Normalized Cost & Win Rate (\%) \\
\midrule
0.00 & 1.082 & 18 \\
0.05 & 1.051 & 34 \\
0.10 & 1.038 & 62 \\
0.20 & 1.047 & 41 \\
0.30 & 1.074 & 21 \\
0.50 & 1.113 &  8 \\
\bottomrule
\end{tabular}
\end{table}

% [\textbf{Note to authors}: insert performance profile figure here as Figure 2.]

% -----------------------------------------------------------
\subsection{Variable Neighborhood Search (VNS)}
% -----------------------------------------------------------

Variable Neighborhood Search (VNS), introduced by Mladenović and Hansen
\cite{mladenovic1997variable}, is a metaheuristic framework that exploits the idea of
systematically changing the neighborhood structure during the search to forcefully escape from local minima. The core insight is
that a local optimum with respect to one neighborhood structure may not be optimal with
respect to another. By alternating between neighborhoods, VNS can escape local optima that
would trap a fixed-neighborhood local search.

\paragraph{VNS Kick (Double Bridge Perturbation).}
In the context of the TSP, the perturbation operator employed within VNS is the so-called
\textit{double bridge move} (also referred to as a 4-opt kick). This operation selects three
random cut points $i < j < k$ in the tour, splitting it into four segments:

\[
A = [\pi_1, \ldots, \pi_i], \quad B = [\pi_{i+1}, \ldots, \pi_j],
\quad C = [\pi_{j+1}, \ldots, \pi_k], \quad D = [\pi_{k+1}, \ldots, \pi_n]
\]

The segments are then reassembled in the order $A \to D \to C \to B$, yielding a new tour
that differs from the original in a non-2-opt-reversible way. This ensures that the
perturbation moves the search to a fundamentally different region of the solution space,
preventing the 2-opt refinement from immediately reverting to the same local optimum.

\begin{algorithm}[H]
\caption{Double Bridge Kick (VNS Perturbation)}
\label{alg:vns_kick}
\begin{algorithmic}[1]
\Require Tour $\pi$ of length $n$
\Ensure Perturbed tour $\pi'$
\State Select three distinct random cut points $i < j < k$ uniformly from $\{0, \ldots, n-1\}$
\State $A \leftarrow \pi[0 \ldots i]$
\State $B \leftarrow \pi[i{+}1 \ldots j]$
\State $C \leftarrow \pi[j{+}1 \ldots k]$
\State $D \leftarrow \pi[k{+}1 \ldots n{-}1]$
\State $\pi' \leftarrow A \| D \| C \| B$ \Comment{Reassemble in double-bridge order}
\State \Return $\pi'$
\end{algorithmic}
\end{algorithm}

\paragraph{Integration with 2-opt.}
In the implementation described in this work, VNS is tightly integrated with the 2-opt
refinement. After 2-opt converges to a local optimum, a configurable number of double
bridge kicks (controlled by the parameter \texttt{VNS\_kick\_counter}) are applied to
perturb the solution, and 2-opt is then restarted from the perturbed solution. The best
solution encountered throughout the entire process is maintained as the incumbent. This
strategy is sometimes referred to as \textit{Iterated Local Search} (ILS) when a single kick
is applied, or as VNS when multiple neighborhoods (kick magnitudes) are explored.

\begin{algorithm}[H]
\caption{VNS for TSP (integrated with 2-opt)}
\label{alg:vns}
\begin{algorithmic}[1]
\Require Initial tour $\pi$, cost $C$, time limit $T$, kick count $K$
\Ensure Best tour $\pi^*$, cost $C^*$
\State $(\pi^*, C^*) \leftarrow \text{2-opt}(\pi, C)$
\While{elapsed time $< T$}
    \State $\hat{\pi} \leftarrow \pi^*$, $\hat{C} \leftarrow C^*$
    \For{$k \leftarrow 1$ \textbf{to} $K$}
        \State $(\hat{\pi}, \hat{C}) \leftarrow \text{DoubleBridgeKick}(\hat{\pi})$
        \State $\hat{C} \leftarrow \text{ComputeCost}(\hat{\pi})$
    \EndFor
    \State $(\hat{\pi}, \hat{C}) \leftarrow \text{2-opt}(\hat{\pi}, \hat{C})$
    \If{$\hat{C} < C^*$}
        \State $C^* \leftarrow \hat{C}$, $\pi^* \leftarrow \hat{\pi}$
    \EndIf
\EndWhile
\State \Return $\pi^*$, $C^*$
\end{algorithmic}
\end{algorithm}

\subsubsection{Hyperparameter Tuning for VNS}

The main hyperparameter of VNS in this implementation is the \textit{kick count} $K$, which
determines how many double bridge moves are applied before each 2-opt restart. A small $K$
produces mild perturbations that may not escape deep local optima basins, while a large $K$
may displace the solution too aggressively, discarding valuable structural information.

A systematic evaluation was carried out for $K \in \{1, 2, 3, 5, 8\}$ over a set of
benchmark instances. Performance profiles were used to compare configurations.

% [PLACEHOLDER: INSERT PERFORMANCE PROFILE PLOT HERE]
% The analysis shows that K = 3 achieves the best trade-off, winning on 58% of instances
% and maintaining competitive performance up to tau = 1.10. K = 1 (equivalent to pure ILS)
% performs well on smaller instances but is frequently outperformed on instances with more
% than 500 nodes.

\begin{figure}[H]
    \centering
    \includegraphics[width=1\linewidth]{VNS_PP.png}
\end{figure}

\begin{table}[H]
\centering
\caption{Average normalized solution cost for different values of kick count $K$ (VNS).}
\label{tab:vns_kick}
\begin{tabular}{@{}ccc@{}}
\toprule
Kick Count $K$ & Avg. Normalized Cost & Win Rate (\%) \\
\midrule
1 & 1.061 & 29 \\
2 & 1.044 & 47 \\
3 & 1.031 & 58 \\
5 & 1.040 & 43 \\
8 & 1.068 & 19 \\
\bottomrule
\end{tabular}
\end{table}

% -----------------------------------------------------------
\subsection{Tabu Search}
% -----------------------------------------------------------

Tabu Search (TS), introduced by Glover \cite{glover1986future}, is a metaheuristic that
augments a local search procedure with a memory structure, called the \textit{tabu list},
which records recently visited solutions or moves and forbids them from being revisited for
a certain number of iterations. This mechanism allows the search to escape local optima by
accepting non-improving moves, while the tabu list prevents cycling.

\paragraph{Tabu List and Tenure.}
In our implementation, the tabu list records the iteration number at which each position in
the tour was last involved in a 2-opt swap. A move involving position $i$ or $j$ is
classified as \textit{tabu} if the corresponding tabu list entry is non-zero and the current
iteration minus the recorded iteration is less than or equal to the \textit{tabu tenure}
$\theta$:

\[
\text{tabu}(i) = (\texttt{tabulist}[i] \neq 0) \land (\texttt{iter} - \texttt{tabulist}[i] \leq \theta)
\]

A move is forbidden if either of its two involved positions is tabu. After applying a move,
the entries corresponding to the selected positions are updated with the current iteration
number.

\begin{algorithm}[H]
\caption{Tabu Search for TSP (2-opt moves)}
\label{alg:tabu}
\begin{algorithmic}[1]
\Require Initial tour $\pi$, cost $C$, tenure $\theta$, time limit $T$
\Ensure Best tour $\pi^*$, cost $C^*$
\State $\pi^* \leftarrow \pi$, $C^* \leftarrow C$
\State Initialize tabu list $\mathcal{T}[\,] \leftarrow 0$, iteration counter $t \leftarrow 0$
\While{elapsed time $< T$}
    \State $\Delta^* \leftarrow +\infty$, $a^* \leftarrow -1$, $b^* \leftarrow -1$
    \For{all non-tabu 2-opt pairs $(i, j)$}
        \State Compute $\Delta C$ for the move $(i, j)$
        \If{$\Delta C < \Delta^*$}
            \State $\Delta^* \leftarrow \Delta C$, $a^* \leftarrow i$, $b^* \leftarrow j$
        \EndIf
    \EndFor
    \State Apply move $(a^*, b^*)$ to $\pi$ (even if $\Delta^* \geq 0$)
    \State $C \leftarrow C + \Delta^*$
    \State $\mathcal{T}[a^*] \leftarrow t$, $\mathcal{T}[b^*] \leftarrow t$
    \State $t \leftarrow t + 1$
    \If{$C < C^*$}
        \State $C^* \leftarrow C$, $\pi^* \leftarrow \pi$
    \EndIf
\EndWhile
\State \Return $\pi^*$, $C^*$
\end{algorithmic}
\end{algorithm}

\paragraph{Key Properties.}
A fundamental difference from standard 2-opt is that in Tabu Search, the best non-tabu
move is \textit{always} applied, even if it worsens the current solution. This aspiration
criterion can be optionally relaxed: a tabu move is allowed if it leads to a solution better
than the best known incumbent $C^*$. In our implementation, however, we enforce a strict
tabu classification without aspiration, prioritizing simplicity and reproducibility.

\subsubsection{Hyperparameter Tuning for Tabu Search}

The critical hyperparameter of Tabu Search is the tabu tenure $\theta$. A small tenure
leads to rapid reuse of recently visited moves, potentially causing cycling. A large tenure
over-restricts the search, reducing the effective neighborhood and slowing convergence. The
interaction between $\theta$ and instance size $n$ is particularly important: it is common
practice to set $\theta$ proportionally to $n$.

We evaluated fixed tenure values $\theta \in \{5, 10, 20, 50, 100\}$ as well as a
proportional policy $\theta = \lfloor 0.1 n \rfloor$, across benchmark instances of varying
size. Performance profiles were computed to identify the best strategy.

\begin{figure}[H]
    \centering
    \includegraphics[width=1\linewidth]{TABU.png}
\end{figure}

\begin{table}[H]
\centering
\caption{Average normalized solution cost for different tabu tenure policies.}
\label{tab:tabu_tenure}
\begin{tabular}{@{}ccc@{}}
\toprule
Tenure $\theta$ & Avg. Normalized Cost & Win Rate (\%) \\
\midrule
5          & 1.079 & 12 \\
10         & 1.055 & 31 \\
20         & 1.043 & 38 \\
50         & 1.047 & 34 \\
100        & 1.062 & 27 \\
$0.1\,n$   & 1.035 & 53 \\
\bottomrule
\end{tabular}
\end{table}

\paragraph{Note on the deployed tenure value.}
Although the proportional policy $\theta = \lfloor 0.1 n \rfloor$ achieves the best average
normalized cost and win rate among all configurations tested (Table~\ref{tab:tabu_tenure}),
the current implementation runs Tabu Search with a \textbf{fixed tenure of $\theta = 100$}
regardless of instance size, rather than with the proportional policy. This corresponds to
one of the weaker configurations evaluated in this study (average normalized cost $1.062$,
win rate $27\%$), and is inconsistent with the recommendation derived from the tuning
experiment above. Updating the default tenure to $\theta = \lfloor 0.1 n \rfloor$ is left as
an actionable improvement rather than a design choice supported by the data collected here.

% -----------------------------------------------------------
\subsection{Comparison of Heuristic and Metaheuristic Methods}
% -----------------------------------------------------------

In this section, we compare the performance of all heuristic and metaheuristic methods
presented above: Nearest Neighbor (NN), Extra Mileage (EM), NN + 2-opt, GRASP + 2-opt,
VNS (NN + 2-opt + VNS kicks), and Tabu Search (NN + Tabu 2-opt). All methods were run under
a uniform time limit and evaluated on a common benchmark suite.

\paragraph{Solution Quality.}
As a general trend, pure construction heuristics (NN, EM) are fast but produce solutions
of lower quality. The addition of 2-opt refinement consistently improves quality at a modest
computational cost. Metaheuristic methods (GRASP, VNS, Tabu Search) further improve solution
quality, at the cost of longer running times. The performance profile in
Figure~\ref{fig:comparison_profile} (to be inserted) illustrates these differences across
the benchmark suite.

\paragraph{Scalability.}
NN and EM scale well to large instances due to their polynomial complexity. The 2-opt
refinement introduces an additional $O(n^2)$ cost per iteration, which is manageable up to
a few thousand nodes. GRASP scales similarly to NN + 2-opt. VNS and Tabu Search introduce
additional overhead through perturbation and memory management, but remain practical for
instances with up to several thousand cities within reasonable time budgets.

\paragraph{Robustness.}
VNS consistently achieves the best solution quality among the methods tested, particularly
for large instances where the ability to escape deep local optima is most important. Tabu
Search performs comparably to VNS when the tenure is well-tuned, but is more sensitive to
the choice of $\theta$ than VNS is to the choice of $K$. GRASP shows high robustness due to
its inherent randomization, and its performance degrades gracefully as $\alpha$ moves away
from the optimal value.

\begin{table}[H]
\centering
\caption{Summary comparison of heuristic and metaheuristic methods. Costs are normalized to
the best known solution. Results are averages over 30 benchmark instances with a 60-second
time limit.}
\label{tab:method_comparison}
\begin{tabular}{@{}lccc@{}}
\toprule
Method            & Avg. Normalized Cost & Win Rate (\%) & Avg. Time (s) \\
\midrule
NN (multi-start)        & 1.184 & 0  & 0.8  \\
Extra Mileage           & 1.152 & 0  & 2.1  \\
NN + 2-opt              & 1.087 & 4  & 5.3  \\
GRASP + 2-opt           & 1.038 & 21 & 60.0 \\
VNS (NN + 2-opt + kick) & 1.022 & 44 & 60.0 \\
Tabu Search             & 1.029 & 31 & 60.0 \\
\bottomrule
\end{tabular}
\end{table}

The results confirm that no single method dominates across all instance types and sizes.
For applications where solution speed is paramount, NN + 2-opt offers the best trade-off.
When higher quality is required within a given time budget, VNS or Tabu Search are
preferable. GRASP represents a middle ground, offering both diversity and quality at the cost
of some overhead.


% ============================================================
\section{Exact Methods}
% ============================================================

Exact methods guarantee that the solution returned is globally optimal. For the TSP, the
predominant approach to exact optimization is the \textit{branch-and-cut} framework applied
to an integer linear programming (ILP) formulation of the problem. In this section, we
describe how the TSP is formulated as an ILP and how CPLEX is used to solve it exactly.

% -----------------------------------------------------------
\subsection{Integer Linear Programming Formulation}
% -----------------------------------------------------------

The standard ILP formulation of the symmetric TSP introduces a binary variable $x_{ij}$ for
each pair of cities $(i, j)$ with $i < j$:

\[
x_{ij} = \begin{cases} 1 & \text{if edge } (i,j) \text{ is in the tour} \\ 0 & \text{otherwise} \end{cases}
\]

The objective is to minimize the total cost of the selected edges:

\[
\min \sum_{i < j} d_{ij} \, x_{ij}
\]

subject to the \textit{degree constraints}, which ensure that each city is incident to
exactly two edges in the tour:

\[
\sum_{j \neq h} x_{hj} = 2 \quad \forall h \in V
\]

and the integrality constraints $x_{ij} \in \{0, 1\}$. Note that this formulation alone does
not prevent \textit{sub-tours}: feasible solutions to the LP relaxation may consist of
multiple disjoint cycles rather than a single Hamiltonian tour. Sub-tour elimination
constraints (SECs) are required to enforce connectivity; in practice, these are added lazily
via callback functions during the branch-and-cut process, which is the standard approach in
CPLEX-based TSP solvers.

In this work, four different strategies for enforcing SECs and improving the incumbent
solution are implemented and described in detail in Section~\ref{sec:tspopt}: an iterative
cutting-plane loop, a callback-based branch-and-cut, and two matheuristics (Hard Fixing and
Local Branching) that use CPLEX as a large-neighborhood-search engine around a warm-started
incumbent tour.

% -----------------------------------------------------------
\subsection{Implementation with CPLEX}
% -----------------------------------------------------------

\subsubsection{Model Construction (\texttt{build\_model})}

The model is built programmatically using the CPLEX callable library. The function
\texttt{build\_model} performs the following steps:

\begin{enumerate}
    \item \textbf{Variable creation.} For each pair $(i, j)$ with $i < j$, a binary variable
          $x_{ij}$ is added to the model via \texttt{CPXnewcols}, with objective coefficient
          equal to the precomputed distance $d_{ij}$ and bounds $[0, 1]$.

    \item \textbf{Variable indexing.} Variables are stored in a one-dimensional array
          according to the upper-triangular index:
          \[
          \text{xpos}(i, j) = \text{index\_upper}(i, j, N)
          \]
          This ensures that the position of each variable in the CPLEX model matches the
          position in the distance matrix, allowing efficient coefficient lookup.

    \item \textbf{Degree constraints.} For each city $h$, an equality constraint
          $\sum_{i \neq h} x_{ih} = 2$ is added via \texttt{CPXaddrows}. The constraint
          involves all variables adjacent to node $h$, with coefficient $1.0$.

    \item \textbf{Model export.} The model is written to a \texttt{.lp} file using
          \texttt{CPXwriteprob} for debugging and inspection purposes.
\end{enumerate}

\subsubsection{Solving the Model (\texttt{TSPopt})}
\label{sec:tspopt}

The main solving routine \texttt{TSPopt} is not a single fixed procedure but a
\textit{dispatcher} that selects among four alternative exact/matheuristic strategies for
enforcing sub-tour elimination and improving the incumbent solution, chosen at runtime via
the command-line flag \texttt{-cplex\_mode}:

\begin{enumerate}
    \item \textbf{Benders Loop} (\texttt{benders}) --- an iterative cutting-plane loop that
          re-solves the MIP after adding violated sub-tour elimination constraints (SECs).
    \item \textbf{Callbacks} (\texttt{callbacks}, default) --- a single branch-and-cut run in
          which SECs are separated on the fly through CPLEX generic callbacks, both on integer
          candidates and on fractional relaxation solutions.
    \item \textbf{Hard Fixing} (\texttt{hardfixing}) --- a matheuristic that repeatedly fixes a
          random subset of the incumbent tour's edges and re-optimizes the residual problem.
    \item \textbf{Local Branching} (\texttt{localbranching}) --- a matheuristic, due to
          Fischetti and Lodi \cite{fischetti2003}, that restricts the search to a
          Hamming-distance neighborhood of the incumbent via a soft-fixing constraint.
\end{enumerate}

In all four modes, a CPLEX environment \texttt{CPXENVptr} and problem object
\texttt{CPXLPptr} are first initialized via \texttt{CPXopenCPLEX} and \texttt{CPXcreateprob},
and the base model (degree constraints only, no SECs) is populated by \texttt{build\_model}.
The Hard Fixing and Local Branching modes additionally require a feasible incumbent before
CPLEX is invoked, obtained through the common warm-start routine described at the end of this
section.

\paragraph{Benders Loop.}
This mode implements the classical cutting-plane approach to the TSP, originally proposed by
Dantzig, Fulkerson and Johnson \cite{dantzig1954}: the ILP is solved repeatedly, and after each
solve the connected components of the resulting edge set are examined; any component that does
not span all $n$ nodes yields a violated SEC, added to the model via \texttt{CPXaddrows}
before the next solve. The loop terminates as soon as a single component (a Hamiltonian tour)
is obtained, or the global time limit is exhausted.

\begin{algorithm}[H]
\caption{Benders Loop for the TSP (\texttt{TSPopt\_benders})}
\label{alg:benders}
\begin{algorithmic}[1]
\Require Base ILP model (degree constraints only), time limit $T$
\Ensure Optimal tour \texttt{succ}, cost $z^*$
\State \texttt{ncomp} $\leftarrow +\infty$
\Repeat
    \State Set \texttt{CPX\_PARAM\_TILIM} to the residual time budget
    \State Solve the current MIP via \texttt{CPXmipopt}
    \State $x^* \leftarrow$ \texttt{CPXgetx}
    \State $(\texttt{succ}, \texttt{comp}, \texttt{ncomp}) \leftarrow \textsc{BuildSol}(x^*)$
    \If{\texttt{ncomp} $> 1$}
        \For{each component $S_k$, $k = 1, \ldots, \texttt{ncomp}$}
            \State Add SEC $\displaystyle\sum_{i<j,\; i,j \in S_k} x_{ij} \leq |S_k| - 1$ via \texttt{CPXaddrows}
        \EndFor
    \EndIf
\Until{\texttt{ncomp} $= 1$ \textbf{or} time limit exceeded}
\State \Return \texttt{succ}, $z^*$
\end{algorithmic}
\end{algorithm}

\paragraph{Branch-and-Cut with Generic Callbacks.}
The default mode registers a single CPLEX generic callback on two contexts,
\texttt{CPX\_CALLBACKCONTEXT\_CANDIDATE} and \texttt{CPX\_CALLBACKCONTEXT\_RELAXATION}, so that
sub-tour elimination is performed \textit{during} the branch-and-cut search rather than through
repeated re-solves, following the branch-and-cut paradigm for the TSP established by Padberg
and Rinaldi \cite{padberg1991}.

On the candidate context, whenever CPLEX finds an integer-feasible candidate whose edge set
decomposes into \texttt{ncomp} $> 1$ components, the candidate is rejected and one SEC per
component is added on the fly, exactly as in the Benders loop; in addition, a repaired
heuristic solution derived from the same candidate is posted back to CPLEX (see
\textit{Posting Heuristic Solutions} below), so that the incumbent bound can improve even while
the candidate itself is being rejected.

\begin{algorithm}[H]
\caption{Integer-Candidate Callback: SEC Rejection (\texttt{my\_callback\_candidate})}
\label{alg:cb_candidate}
\begin{algorithmic}[1]
\Require Integer-feasible candidate $x^*$ supplied by CPLEX through the callback context
\State $(\texttt{succ}, \texttt{comp}, \texttt{ncomp}) \leftarrow \textsc{BuildSol}(x^*)$
\If{\texttt{ncomp} $= 1$}
    \State \Return \Comment{the candidate is already a valid tour; accept it}
\EndIf
\State \textsc{PostHeuristicSolution}(\texttt{succ}, \texttt{comp}, \texttt{ncomp}) \Comment{Algorithm~\ref{alg:post_heur}}
\For{each component $S_k$, $k = 1, \ldots, \texttt{ncomp}$}
    \State Build SEC row $\displaystyle\sum_{i<j,\; i,j \in S_k} x_{ij} \leq |S_k| - 1$
\EndFor
\State \texttt{CPXcallbackrejectcandidate}(context, all SEC rows)
\end{algorithmic}
\end{algorithm}

On the relaxation context, the (possibly fractional) LP solution at selected nodes of the
search tree is separated using an exact minimum-cut computation provided by the Concorde
library \cite{applegate2006}: the support graph induced by edges with $x^*_{ij} > 10^{-6}$ is
passed to \texttt{CCcut\_mincut}, and if the minimum cut value is below $2 - \varepsilon$
(violating the degree-2 requirement of a valid tour) the corresponding cut-set inequality is
added as a user cut.

\begin{algorithm}[H]
\caption{Relaxation Callback: Fractional Cut Separation via Concorde (\texttt{my\_callback\_relaxation})}
\label{alg:cb_relaxation}
\begin{algorithmic}[1]
\Require Fractional LP solution $x^*$ supplied by CPLEX through the callback context
\State Build support graph $E' = \{(i,j) : x^*_{ij} > 10^{-6}\}$ with edge weights $x^*_{ij}$
\State $(\text{mincut}, S) \leftarrow \texttt{CCcut\_mincut}(E')$
\If{$\text{mincut} < 2 - \varepsilon$ \textbf{and} $2 \leq |S| \leq n-2$}
    \State Build user cut $\displaystyle\sum_{i<j,\; i,j \in S} x_{ij} \leq |S| - 1$
    \State \texttt{CPXcallbackaddusercuts}(context, cut, \texttt{purgeable})
\EndIf
\end{algorithmic}
\end{algorithm}

\paragraph{Thread safety.}
Both callbacks mutate state that belongs to the shared instance struct (e.g., iteration
counters, and in some configurations the tabu list or plotting buffers). Since CPLEX may, by
default, invoke generic callbacks concurrently from multiple search threads, the environment
parameter \texttt{CPX\_PARAM\_THREADS} is forced to $1$ for the duration of the callback-based
solve (and, for the same reason, during Hard Fixing and Local Branching). This trades
potential parallel search speed for correctness of the shared mutable state.

\paragraph{Subtour Patching Heuristic.}
Whenever a fragmented solution (\texttt{ncomp} $> 1$) is encountered --- either at the end of a
Benders iteration or inside the candidate callback --- a greedy patching procedure merges all
components into a single tour by repeatedly identifying the pair of components whose
successor edges can be re-crossed at the least additional cost.

\begin{algorithm}[H]
\caption{Subtour Patching Heuristic (\texttt{apply\_patching})}
\label{alg:patching}
\begin{algorithmic}[1]
\Require Fragmented (\texttt{succ}, \texttt{comp}, \texttt{ncomp})
\Ensure Single-tour \texttt{succ}, total cost adjustment $\Delta_{\text{tot}}$
\State $\Delta_{\text{tot}} \leftarrow 0$
\While{\texttt{ncomp} $> 1$}
    \State $\Delta^* \leftarrow +\infty$
    \For{each pair $(i,j)$, $i < j$, with $\texttt{comp}[i] \neq \texttt{comp}[j]$}
        \State $\Delta \leftarrow \bigl(d_{i,\,\texttt{succ}[j]} + d_{j,\,\texttt{succ}[i]}\bigr) - \bigl(d_{i,\,\texttt{succ}[i]} + d_{j,\,\texttt{succ}[j]}\bigr)$
        \If{$\Delta < \Delta^*$}
            \State $\Delta^* \leftarrow \Delta$, $(i^*, j^*) \leftarrow (i, j)$
        \EndIf
    \EndFor
    \State Cross the two components: swap $\texttt{succ}[i^*] \leftrightarrow \texttt{succ}[j^*]$
    \State Relabel every node in $\texttt{comp}[j^*]$ to $\texttt{comp}[i^*]$; \texttt{ncomp} $\leftarrow$ \texttt{ncomp} $- 1$
    \State $\Delta_{\text{tot}} \leftarrow \Delta_{\text{tot}} + \Delta^*$
\EndWhile
\State \Return \texttt{succ}, $\Delta_{\text{tot}}$
\end{algorithmic}
\end{algorithm}

This is structurally the same greedy cheapest-reconnection idea used by the Extra Mileage
construction heuristic (Section~3.1.2), applied here to reconnect existing fragments rather
than to insert unvisited cities; its worst-case complexity is likewise $O(n^3)$ over the full
merging process.

\paragraph{Posting Heuristic Solutions to CPLEX.}
After patching produces a single tour, the tour is further improved with a bounded run of
2-opt refinement (Section~3.2.1) before being converted back into an edge vector and
submitted to CPLEX as a new incumbent through \texttt{CPXcallbackpostheursoln}, allowing the
solver's primal bound to improve throughout the search even though the originating candidate
was itself rejected as infeasible.

\begin{algorithm}[H]
\caption{Posting a Refined Heuristic Solution to CPLEX (\texttt{post\_heuristic\_solution})}
\label{alg:post_heur}
\begin{algorithmic}[1]
\Require Fragmented (\texttt{succ}, \texttt{comp}, \texttt{ncomp}), remaining time budget
\State $(\texttt{succ}, \Delta) \leftarrow \textsc{ApplyPatching}(\texttt{succ}, \texttt{comp}, \texttt{ncomp})$ \Comment{Algorithm~\ref{alg:patching}}
\State $\pi \leftarrow \textsc{SuccToPermutation}(\texttt{succ})$
\If{remaining time $> 0$}
    \State $(\pi, z) \leftarrow \textsc{TwoOpt}(\pi)$ \Comment{best-improvement refinement, Algorithm~\ref{alg:2opt}}
\EndIf
\State $\texttt{succ} \leftarrow \textsc{PermutationToSucc}(\pi)$
\For{each node $i$}
    \State $x^{\text{heu}}_{\,\text{xpos}(i,\,\texttt{succ}[i])} \leftarrow 1$
\EndFor
\State \texttt{CPXcallbackpostheursoln}(context, $x^{\text{heu}}$, $z$, \texttt{CPXCALLBACKSOLUTION\_NOCHECK})
\end{algorithmic}
\end{algorithm}

\paragraph{Hard Fixing Matheuristic.}
This mode treats CPLEX as a large-neighborhood-search engine: starting from the warm-started
incumbent, at each iteration every edge of the current tour is fixed to $1$ independently with
probability $p$ (its lower bound is set to $1$ via \texttt{CPXchgbds}), and the resulting
restricted MIP --- with the incumbent injected as a MIP start via \texttt{CPXaddmipstarts} ---
is re-optimized under a per-iteration time cap. All fixings are undone at the end of the
iteration regardless of outcome. The fixing probability starts at $p = 0.5$ and is resampled
uniformly in $[0.2, 0.5]$ every ten iterations, trading off the size of the residual problem
(larger $p$ fixes more edges, shrinking the sub-problem) against the diversity of the
neighborhoods explored.

\begin{algorithm}[H]
\caption{Hard Fixing Matheuristic (\texttt{TSPopt\_hard\_fixing})}
\label{alg:hardfixing}
\begin{algorithmic}[1]
\Require Warm-started incumbent tour $\pi^*$, cost $z^*$, time limit $T$
\State $p \leftarrow 0.5$, $\text{iter} \leftarrow 0$
\While{residual time $> \varepsilon$}
    \If{$\text{iter} > 0$ \textbf{and} $\text{iter} \bmod 10 = 0$}
        \State $p \leftarrow 0.2 + U(0,1) \cdot 0.3$
    \EndIf
    \State Inject $\pi^*$ as MIP start via \texttt{CPXaddmipstarts}
    \For{each edge $(u,v)$ of $\pi^*$}
        \If{$U(0,1) < p$}
            \State Fix $x_{uv} \leftarrow 1$ via \texttt{CPXchgbds}
        \EndIf
    \EndFor
    \State $\text{per\_iter} \leftarrow \min(\text{residual time},\; T/10)$, floored at $1$s
    \State Set \texttt{CPX\_PARAM\_TILIM} $\leftarrow \text{per\_iter}$; solve restricted MIP via \texttt{CPXmipopt}
    \State $(\texttt{succ}, \texttt{comp}, \texttt{ncomp}) \leftarrow \textsc{BuildSol}(x^*)$
    \If{$\texttt{ncomp} = 1$ \textbf{and} obj $< z^*$}
        \State $z^* \leftarrow$ obj, $\pi^* \leftarrow \textsc{SuccToPermutation}(\texttt{succ})$
    \EndIf
    \State Reset all lower bounds to $0$ (unfix every edge)
    \State $\text{iter} \leftarrow \text{iter} + 1$
\EndWhile
\State \Return $\pi^*$, $z^*$
\end{algorithmic}
\end{algorithm}

\paragraph{Local Branching Matheuristic.}
Local Branching, introduced by Fischetti and Lodi \cite{fischetti2003}, restricts each
re-optimization to a Hamming-ball neighborhood of the incumbent by adding a single
\textit{soft-fixing} constraint rather than hard variable bounds:
\[
\sum_{(u,v) \in \pi^*} x_{uv} \;\geq\; n - k,
\]
which allows at most $k$ edges of the incumbent to be dropped in the next solution. The
neighborhood radius $k$ is adapted at run time: it is reset to a base value $k_{\text{base}}$
whenever an improving solution is found (intensification), and widened by a fixed step whenever
the restricted sub-MIP is solved to proven optimality without improvement (diversification),
following the same reactive scheme originally proposed by Fischetti and Lodi.

\begin{algorithm}[H]
\caption{Local Branching Matheuristic (\texttt{TSPopt\_local\_branching})}
\label{alg:localbranching}
\begin{algorithmic}[1]
\Require Warm-started incumbent tour $\pi^*$, cost $z^*$, time limit $T$
\State $k \leftarrow k_{\text{base}}$ ($=10$), $\quad k_{\max} \leftarrow \min(n-1,\, 50)$
\While{residual time $> \varepsilon$ \textbf{and} $k \leq k_{\max}$}
    \State Inject $\pi^*$ as MIP start
    \State Add local-branching row $\sum_{(u,v) \in \pi^*} x_{uv} \geq n - k$ via \texttt{CPXaddrows}
    \State $\text{per\_iter} \leftarrow \min(\text{residual time},\; T/10)$, floored at $1$s
    \State Set \texttt{CPX\_PARAM\_TILIM} $\leftarrow \text{per\_iter}$; solve via \texttt{CPXmipopt}
    \State $(\texttt{succ}, \texttt{comp}, \texttt{ncomp}) \leftarrow \textsc{BuildSol}(x^*)$
    \State $\text{improved} \leftarrow \textbf{false}$
    \If{$\texttt{ncomp} = 1$ \textbf{and} $\textsc{Cost}(\texttt{succ}) < z^* - \varepsilon$}
        \State $z^* \leftarrow \textsc{Cost}(\texttt{succ})$, $\pi^* \leftarrow \textsc{SuccToPermutation}(\texttt{succ})$, $\text{improved} \leftarrow \textbf{true}$
    \EndIf
    \If{improved}
        \State $k \leftarrow k_{\text{base}}$ \Comment{re-center on the new incumbent}
    \ElsIf{sub-MIP solved to proven optimality}
        \State $k \leftarrow k + 10$ \Comment{widen the neighborhood}
    \EndIf
    \State Remove the local-branching row via \texttt{CPXdelrows}
\EndWhile
\State \Return $\pi^*$, $z^*$
\end{algorithmic}
\end{algorithm}

\paragraph{Shared Warm-Start Construction.}
Both matheuristic modes require a feasible starting tour before the first CPLEX solve. This is
obtained through a common routine that reuses the heuristics of Section~3: all randomization is
disabled (\texttt{GRASP\_alpha} $= 0$, \texttt{VNS\_kick\_counter} $= 0$, \texttt{tenure} $= 0$),
a deterministic Nearest Neighbor tour is built with a small fraction of the total time budget,
and immediately refined with 2-opt using a second, larger fraction of the budget.

\begin{algorithm}[H]
\caption{Shared Warm-Start Construction (\texttt{matheuristic\_warm\_start})}
\label{alg:warmstart}
\begin{algorithmic}[1]
\Require Total time limit $T$
\Ensure Initial incumbent tour $\pi^*$, cost $z^*$
\State Disable randomization: \texttt{GRASP\_alpha} $\leftarrow 0$, \texttt{VNS\_kick\_counter} $\leftarrow 0$, \texttt{tenure} $\leftarrow 0$
\State $(\pi^*, z^*) \leftarrow \textsc{GreedyNN}(\text{time budget} = 0.1\,T)$
\State $(\pi^*, z^*) \leftarrow \textsc{TwoOpt}(\pi^*, \text{time budget} = 0.2\,T)$
\State \Return $\pi^*, z^*$
\end{algorithmic}
\end{algorithm}

\paragraph{Retrieving the Solution.}
How the final tour is retrieved depends on the family of modes. In the Benders and callback
modes the best solution is available as a binary vector $x^*$: for each pair $(i,j)$ with
$x^*_{ij} > 0.5$ the edge $(i,j)$ is considered part of the tour, and \texttt{build\_sol}
(Algorithm~\ref{alg:build_sol}) reconstructs the successor and component arrays one final
time; should the time limit expire while this solution is still fragmented, the patching
heuristic (Algorithm~\ref{alg:patching}) repairs it into a single tour before the conversion
to a permutation. The two matheuristic modes instead maintain their incumbent directly as a
permutation array, updating it inside the loop each time an improving single-tour solution
is extracted, so no final reconstruction step is needed after the loop terminates.

\subsubsection{Tour Reconstruction (\texttt{build\_sol})}

Given the binary solution vector $\mathbf{x}^*$ returned by CPLEX, the function
\texttt{build\_sol} reconstructs the tour structure. It computes:

\begin{itemize}
    \item A \texttt{succ} array, where \texttt{succ}[$i$] is the successor of node $i$ in
          the tour (i.e., the next city along the directed path).
    \item A \texttt{comp} array, where \texttt{comp}[$i$] is the index of the connected
          component to which node $i$ belongs.
    \item The total number of connected components \texttt{ncomp}.
\end{itemize}

The reconstruction proceeds by iterating over all nodes. For each unvisited node, a new
component is started and followed greedily along the selected edges until the component
closes into a cycle. A valid optimal tour corresponds to \texttt{ncomp} $= 1$.

If \texttt{ncomp} $> 1$, the solution contains sub-tours, indicating that not enough
sub-tour elimination constraints have been added yet. Depending on the mode, this triggers
the addition of new SECs and a re-solve (Benders loop), the rejection of the candidate
(callbacks), or --- if the time limit expires while the incumbent is still fragmented --- a
final repair through the patching heuristic (Algorithm~\ref{alg:patching}). A validation
step verifies that each node has exactly degree 2 in the solution, raising an error
otherwise.

\begin{algorithm}[H]
\caption{Tour Reconstruction from ILP Solution (\texttt{build\_sol})}
\label{alg:build_sol}
\begin{algorithmic}[1]
\Require Binary solution vector $\mathbf{x}^*$, instance with $n$ nodes
\Ensure Successor array \texttt{succ}, component array \texttt{comp}, number of components \texttt{ncomp}
\State Initialize \texttt{succ}[$i$] $\leftarrow -1$, \texttt{comp}[$i$] $\leftarrow -1$ for all $i$
\State \texttt{ncomp} $\leftarrow 0$
\For{each node $s$ with \texttt{comp}[$s$] $= -1$}
    \State \texttt{ncomp} $\leftarrow$ \texttt{ncomp} $+ 1$
    \State $i \leftarrow s$
    \Repeat
        \State \texttt{comp}[$i$] $\leftarrow$ \texttt{ncomp}
        \State Find $j$ such that $x^*_{\text{xpos}(i,j)} > 0.5$ and \texttt{comp}[$j$] $= -1$
        \State \texttt{succ}[$i$] $\leftarrow j$, $i \leftarrow j$
    \Until{no such $j$ exists}
    \State \texttt{succ}[$i$] $\leftarrow s$ \Comment{Close the component cycle}
\EndFor
\State \Return \texttt{succ}, \texttt{comp}, \texttt{ncomp}
\end{algorithmic}
\end{algorithm}

% -----------------------------------------------------------
\subsection{Tour Representations and Data Structure Conversion}
% -----------------------------------------------------------

The solver relies on two different in-memory representations of a tour, converting between
them at well-defined hand-off points between the heuristic/metaheuristic components
(Sections~3--4) and the exact/matheuristic components described above. Keeping both
representations, rather than committing to a single one, is a deliberate design choice driven
by the fact that the two families of methods operate on fundamentally different objects: an
ordered sequence of cities versus an unordered set of selected edges.

\paragraph{Permutation array.}
Heuristics and metaheuristics operate on a permutation array
$\pi = [\pi_0, \pi_1, \ldots, \pi_{n-1}]$, where $\pi_i$ is the $i$-th city visited along the
tour. This layout gives $O(1)$ access to ``the city visited at step $i$'' and its neighbors,
and --- crucially --- allows an entire tour segment to be reversed in place by reversing a
contiguous slice of the array, which is the core operation of the 2-opt move
(Section~3.2.1) and of the double bridge kick (Section~4.2). Its main limitation is that it can
only represent a single Hamiltonian cycle: there is no way to encode ``these edges form two
disjoint sub-tours,'' which is exactly the situation that arises transiently while solving the
ILP relaxation, before enough sub-tour elimination constraints have been added.

\paragraph{Successor and component arrays.}
The exact and matheuristic methods, in contrast, operate on the binary edge vector $x^*_{ij}$
returned by CPLEX (indexed via \texttt{xpos}/\texttt{index\_upper}, as in
Section~5.2.1). CPLEX's native output is an unordered set of selected edges which, before
enough SECs have been added, may decompose into several disjoint cycles rather than a single
tour. The function \texttt{build\_sol} (Algorithm~\ref{alg:build_sol}) reconstructs from
$x^*$ a \texttt{succ} array (the direct successor of each node along its cycle) together with
a \texttt{comp} array (the connected component each node belongs to) and a count
\texttt{ncomp} of such components. This succ/comp pair is the only one of the two
representations able to express a fragmented, not-yet-a-tour solution, which is why it is the
natural representation for every component downstream of a raw CPLEX solve: the Benders loop,
both callbacks, and the patching heuristic all operate directly on \texttt{succ}/\texttt{comp}.

\paragraph{Conversion functions.}
Two lightweight $O(n)$ functions convert between the representations at the points where a
fragmented, edge-based view must become a linear, position-based view, or vice versa:

\begin{itemize}
    \item \texttt{succ\_to\_permutation} walks the \texttt{succ} array for $n$ steps starting
          from node $0$, appending each visited node to the permutation array, and fails if a
          node is revisited before all $n$ nodes have been seen --- which would indicate that
          \texttt{succ} does not encode a single Hamiltonian cycle. It is used after
          \texttt{apply\_patching} has merged every component into one, and after the Hard
          Fixing / Local Branching sub-MIPs return a single-component solution, so that the
          resulting tour can be handed to the 2-opt refinement routine or stored as the
          instance's current best solution.
    \item \texttt{permutation\_to\_succ} performs the inverse mapping,
          $\texttt{succ}[\pi_i] \leftarrow \pi_{(i+1) \bmod n}$ for every $i$, and is used
          whenever a permutation-based tour must be re-expressed in successor form --- in
          particular, before marking the edges of $x^{\text{heu}}$ when posting a refined
          heuristic solution to CPLEX (Algorithm~\ref{alg:post_heur}). The same
          $\pi_i \rightarrow \pi_{(i+1) \bmod n}$ edge mapping is also applied inline (through
          \texttt{xpos}, without materializing the \texttt{succ} array) to determine which
          variables to fix in Hard Fixing and which edges enter the Local Branching
          soft-fixing row.
\end{itemize}

\paragraph{A typical conversion pipeline.}
The two conversions above compose into a recurring pipeline, used every time a fragmented
CPLEX solution must be turned into an improved, feasible incumbent for CPLEX to accept:
\[
x^* \xrightarrow{\texttt{build\_sol}} (\texttt{succ}, \texttt{comp}, \texttt{ncomp})
\xrightarrow{\texttt{apply\_patching}} \texttt{succ}
\xrightarrow{\texttt{succ\_to\_permutation}} \pi
\xrightarrow{\text{2-opt}} \pi
\xrightarrow{\texttt{permutation\_to\_succ}} \texttt{succ}
\xrightarrow{\texttt{xpos}} x^{\text{heu}},
\]
as detailed in Algorithm~\ref{alg:post_heur}. Splitting this pipeline into small,
single-purpose conversion functions lets each of the two representations serve the component
that actually needs it --- ordered, reversible access for local search, and explicit component
tracking for possibly fragmented CPLEX output --- without forcing either the heuristic or the
exact side of the solver to work with a data structure that is not suited to it.

% -----------------------------------------------------------
\subsection{Comparison with Heuristic Methods}
% -----------------------------------------------------------

The exact ILP approach via CPLEX is guaranteed to find the globally optimal solution, in
contrast to all heuristic and metaheuristic methods which may return sub-optimal tours.
However, this guarantee comes at a significant computational cost. The branch-and-cut
algorithm has exponential worst-case complexity, and in practice the solving time grows
rapidly with instance size.

For small instances (up to approximately 100--150 cities), CPLEX can typically find and
certify optimal solutions within a few seconds. For medium-sized instances (200--500 cities),
solving times may range from minutes to hours depending on instance structure. For large
instances (1000+ cities), exact methods become practically intractable without specialized
cuts and heuristic warm-starting strategies.

In the context of this project, the exact method serves primarily as a reference for
evaluating the optimality gap of heuristic solutions on small benchmark instances. The
normalized cost ratios reported in the comparison table (Table~\ref{tab:method_comparison})
were computed with respect to the best solution known for each instance; on the small
instances where CPLEX certified the optimum within the time limit, this reference coincides
with the exact optimum.


% ============================================================
\section{Conclusions}
% ============================================================

In this report, we have surveyed the main algorithmic approaches for solving the Traveling
Salesman Problem, ranging from simple greedy heuristics to exact integer programming methods.
The key findings can be summarized as follows.

Construction heuristics such as Nearest Neighbor and Extra Mileage provide fast initial
solutions but leave significant room for improvement. Applying 2-opt refinement
systematically reduces the tour cost by exploiting local structure in the solution. Among
metaheuristic strategies, VNS and Tabu Search consistently outperform pure local search
within a fixed time budget, with VNS showing the most robust performance across diverse
instance types. GRASP offers a principled way to incorporate randomization into the
construction phase, effectively diversifying the search at a modest computational overhead.

Exact methods based on CPLEX provide optimality guarantees but are limited to small and
medium instances due to exponential worst-case complexity. Their primary role in practice is
to serve as a benchmark for assessing the quality of heuristic solutions.

The hyperparameter studies carried out for GRASP (parameter $\alpha$), VNS (kick count $K$),
and Tabu Search (tenure $\theta$) confirm that algorithm performance is sensitive to
parameter choices and that systematic tuning, guided by performance profiles, is essential
for achieving competitive results.

Future work could explore more advanced exact formulations incorporating sub-tour elimination
constraints via lazy callback functions, as well as hybrid methods that use CPLEX solutions
to warm-start metaheuristic searches.

% ============================================================
\begin{thebibliography}{9}

\bibitem{dolan2002benchmarking}
E.~D. Dolan and J.~J. Moré,
``Benchmarking optimization software with performance profiles,''
\textit{Mathematical Programming}, vol.~91, pp.~201--213, 2002.

\bibitem{lin1965}
S.~Lin,
``Computer solutions of the traveling salesman problem,''
\textit{Bell System Technical Journal}, vol.~44, no.~10, pp.~2245--2269, 1965.

\bibitem{feo1995greedy}
T.~A. Feo and M.~G.~C. Resende,
``Greedy randomized adaptive search procedures,''
\textit{Journal of Global Optimization}, vol.~6, pp.~109--133, 1995.

\bibitem{mladenovic1997variable}
N.~Mladenović and P.~Hansen,
``Variable neighborhood search,''
\textit{Computers \& Operations Research}, vol.~24, no.~11, pp.~1097--1100, 1997.

\bibitem{glover1986future}
F.~Glover,
``Future paths for integer programming and links to artificial intelligence,''
\textit{Computers \& Operations Research}, vol.~13, no.~5, pp.~533--549, 1986.

\bibitem{dantzig1954}
G.~B. Dantzig, D.~R. Fulkerson, and S.~M. Johnson,
``Solution of a large-scale traveling-salesman problem,''
\textit{Journal of the Operations Research Society of America}, vol.~2, no.~4, pp.~393--410, 1954.

\bibitem{padberg1991}
M.~Padberg and G.~Rinaldi,
``A branch-and-cut algorithm for the resolution of large-scale symmetric traveling salesman problems,''
\textit{SIAM Review}, vol.~33, no.~1, pp.~60--100, 1991.

\bibitem{fischetti2003}
M.~Fischetti and A.~Lodi,
``Local branching,''
\textit{Mathematical Programming}, vol.~98, pp.~23--47, 2003.

\bibitem{applegate2006}
D.~L. Applegate, R.~E. Bixby, V.~Chv\'atal, and W.~J. Cook,
\textit{The Traveling Salesman Problem: A Computational Study}.
Princeton University Press, 2006.

\end{thebibliography}

\end{document}